\documentclass[10pt]{article}
% Shared LaTeX style for MIT 18.06 exam revision notes.
% Compile topic files with Tectonic, XeLaTeX, or LuaLaTeX.

\usepackage[a4paper,margin=0.72in]{geometry}
\usepackage{fontspec}
\usepackage{amsmath,amssymb,mathtools}
\usepackage{booktabs,array,tabularx}
\usepackage{enumitem}
\usepackage{titlesec}
\usepackage{fancyhdr}
\usepackage{microtype}
\usepackage{xcolor}

\setmainfont{Times New Roman}
\setsansfont{Arial}

\setlength{\parindent}{0pt}
\setlength{\parskip}{3.6pt}
\setlength{\abovedisplayskip}{6pt}
\setlength{\belowdisplayskip}{6pt}
\setlength{\abovedisplayshortskip}{4pt}
\setlength{\belowdisplayshortskip}{4pt}
\renewcommand{\arraystretch}{1.08}
\setlength{\tabcolsep}{4.5pt}

\titleformat{\section}
  {\normalfont\large\bfseries}
  {\thesection}
  {0.55em}
  {}
\titlespacing*{\section}{0pt}{10pt}{3pt}

\titleformat{\subsection}
  {\normalfont\normalsize\bfseries}
  {\thesubsection}
  {0.5em}
  {}
\titlespacing*{\subsection}{0pt}{7pt}{2pt}

\setlist[itemize]{leftmargin=1.35em,itemsep=1pt,topsep=2pt,parsep=0pt}
\setlist[enumerate]{leftmargin=1.55em,itemsep=1pt,topsep=2pt,parsep=0pt}

\pagestyle{fancy}
\fancyhf{}
\fancyfoot[C]{\scriptsize\color{black!45}MIT 18.06 Linear Algebra Exam Notes --- \thepage}
\renewcommand{\headrulewidth}{0pt}
\renewcommand{\footrulewidth}{0pt}

\newcommand{\notesubtitle}[1]{%
  {\centering\itshape #1\par}
  \vspace{2pt}
}

\newcommand{\source}[1]{%
  {\centering\small #1\par}
  \vspace{6pt}
}

\newcommand{\labelnote}[2]{%
  \par\smallskip
  \noindent\textbf{#1.}\quad #2
  \par\smallskip
}

\newcommand{\tightlabel}[1]{%
  \par\smallskip
  \noindent\textbf{#1.}\quad
}

\newcolumntype{Y}{>{\raggedright\arraybackslash}X}
\newcolumntype{L}[1]{>{\raggedright\arraybackslash}p{#1}}

\newenvironment{notetable}
  {\small\begin{tabularx}{\linewidth}}
  {\end{tabularx}}


\begin{document}

\begin{center}
{\LARGE 2.6--2.7 LU, LDU, Permutations, and Transposes\par}
\end{center}
\notesubtitle{Concise exam revision notes for elimination factorizations and matrix operation rules.}
\source{Based on handwritten MIT 18.06 revision notes.}

\section{Core Idea}

\labelnote{Core idea}{Elimination can be stored as a factorization. Instead of redoing row operations, write the matrix as products of simple matrices.}

Without row exchanges, elimination gives
\[
A=LU.
\]
If the pivots are separated into a diagonal matrix \(D\), then
\[
A=LDU.
\]
Here \(L\) stores elimination multipliers, \(D\) stores pivots, and \(U\) stores the final upper triangular structure.

\section{Solving With LU}

To solve
\[
Ax=b,
\qquad A=LU,
\]
do two triangular solves:
\[
Lc=b,
\qquad
Ux=c.
\]

\begin{center}
\small
\begin{tabularx}{\linewidth}{@{}L{0.25\linewidth}Y Y@{}}
\toprule
Step & Equation & Meaning \\
\midrule
Forward substitution & \(Lc=b\) & Apply stored elimination. \\
Back substitution & \(Ux=c\) & Solve the upper triangular system. \\
\bottomrule
\end{tabularx}
\end{center}

\labelnote{Exam memory}{LU is useful because solving triangular systems is faster than starting elimination again.}

\section{Permutations}

If a row exchange is needed, use a permutation matrix:
\[
PA=LU.
\]
Permutation matrices reorder rows and satisfy
\[
P^{-1}=P^{T},
\qquad
P^{T}P=I.
\]

\begin{center}
\small
\begin{tabularx}{\linewidth}{@{}L{0.30\linewidth}Y@{}}
\toprule
Fact & Meaning \\
\midrule
\(P^{-1}=P^{T}\) & Undoing a permutation is another permutation. \\
\(P^{T}P=I\) & Permutations preserve lengths and dot products. \\
\(PA\) & Reorders rows of \(A\). \\
\(AP\) & Reorders columns of \(A\). \\
\bottomrule
\end{tabularx}
\end{center}

\section{Transpose and Inverse Rules}

Order reverses when transposing or inverting products:
\[
(AB)^{T}=B^{T}A^{T},
\qquad
(AB)^{-1}=B^{-1}A^{-1}.
\]
For an inverse transpose,
\[
(A^{-1})^{T}=(A^{T})^{-1}.
\]

\labelnote{Common mistake}{Do not transpose a product in the same order. The order must reverse.}

\section{Symmetric Matrices}

A matrix is symmetric when
\[
A^{T}=A.
\]
The product \(A^{T}A\) is always symmetric:
\[
(A^{T}A)^{T}=A^{T}A.
\]
The dot-product identity behind this is
\[
(Ax)^{T}y=x^{T}A^{T}y.
\]

\section{Block Transposes}

Transpose works block by block and also swaps block positions:
\[
\begin{bmatrix}
A&B\\
C&D
\end{bmatrix}^{T}
=
\begin{bmatrix}
A^{T}&C^{T}\\
B^{T}&D^{T}
\end{bmatrix}.
\]
This is the same rule as ordinary transposes: rows become columns.

\section{Banded Matrices}

Banded matrices have nonzeros near the diagonal, so elimination can be cheaper when the band structure is preserved. A common example is tridiagonal structure:
\[
\begin{bmatrix}
*&*&0&0\\
*&*&*&0\\
0&*&*&*\\
0&0&*&*
\end{bmatrix}.
\]
The exam habit is to use the structure instead of treating every entry as dense.

\section{Common Mistakes}

\begin{center}
\small
\begin{tabularx}{\linewidth}{@{}L{0.40\linewidth}Y@{}}
\toprule
Mistake & Correct exam habit \\
\midrule
Forgetting row exchanges in LU. & Use \(PA=LU\) when pivots require swaps. \\
Solving \(LUx=b\) in one jump. & First solve \(Lc=b\), then \(Ux=c\). \\
Writing \((AB)^{T}=A^{T}B^{T}\). & Reverse order: \(B^{T}A^{T}\). \\
Assuming \(A^{T}A=A A^{T}\). & Both are symmetric, but usually different. \\
Ignoring zeros in a banded matrix. & Use the sparsity pattern to simplify work. \\
\bottomrule
\end{tabularx}
\end{center}

\section{Final Exam Checklist}

\tightlabel{Checklist}
\begin{enumerate}
  \item Check whether row exchanges are needed.
  \item Use \(A=LU\) if no row exchanges occur.
  \item Use \(PA=LU\) if row exchanges occur.
  \item Solve \(Lc=b\) by forward substitution.
  \item Solve \(Ux=c\) by back substitution.
  \item Reverse order when transposing or inverting products.
  \item Check symmetry with \(A^{T}=A\).
  \item Use banded or triangular structure when present.
\end{enumerate}

\end{document}
